% Shared module: exponential-polynomials.
\subsection{Exponential-polynomials: the solutions that remember too little}\label{sec:ep}
\begin{definition}[Exponential-polynomials]\label{def:ep}
$\EP$ is the set of finite sums $\sum_{j}p_j(t,x)\,e^{\tau_jt+\xi_jx}$ with polynomials $p_j$ and
exponents $(\tau_j,\xi_j)\in\C^2$. $\EP(\Omega)$ denotes their restrictions to $\Omega$.
\end{definition}

For a function $w$ let $\ring w=\spn\{\partial_t^{\,j}\partial_x^{\,k}w:\ j,k\ge0\}$ be the space of all its
derivatives. The map $Q\mapsto Q(\partial)w$ is onto $\ring w$ with kernel $\Ann(w)$, so
$\ring w\cong\ring/\Ann(w)$ as vector spaces.

\begin{lemma}[Finite-dimension criterion]\label{lem:ep}
Let $\Omega$ be connected. Then $w\in\EP(\Omega)$ if and only if $\dim_\C\ring w<\infty$. In that case
$\Var(\Ann w)$ is the finite set of exponents $\{(\tau_j,\xi_j)\}$ that occur in $w$.
\end{lemma}

\begin{proof}[Proof idea]
We show that $w$ is an exponential-polynomial if and only if its derivatives span a finite-dimensional
space. ($\Rightarrow$) For one term $w=p(t,x)\,e^{at+bx}$,
\[
  \partial_tw=(p_t+a\,p)\,e^{at+bx},\qquad \partial_xw=(p_x+b\,p)\,e^{at+bx}.
\]
Differentiation never raises the degree of $p$, so every derivative of $w$ lies in the finite-dimensional space
of polynomials of degree at most $\deg p$ times $e^{at+bx}$; a finite sum of such terms stays in a finite sum of
such spaces. ($\Leftarrow$) If $\ring w$ is finite-dimensional, $\partial_t$ and $\partial_x$ map it to itself and act
on it as two commuting matrices $A$ and $B$. Their minimal polynomials $m_A$ and $m_B$ give two
constant-coefficient differential equations,
\[
  m_A(\partial_t)\,w=0,\qquad m_B(\partial_x)\,w=0,
\]
one in each variable. On every rectangle inside $\Omega$ their common solutions are exponential-polynomials, and
connectedness extends this representation to all of $\Omega$. The joint eigenvalues $(a,b)$ of $A$ and $B$ are
the exponents that occur in $w$. Details are in \cref{app:proofs}.
\end{proof}

\inbrief{Deciding whether $w\in\EP$ is a finite-dimensional linear-algebra question: compute the rank
of the derivatives of $w$. The joint eigenvalues of the two matrices are the exponents, that is, the
points of $\C^2$ at which $w$ lives. This is the computational core of everything that follows: the
matrix pencil of \cref{sec:alg-pencil} is exactly this eigenvalue computation, done for $\partial_t$
one wavenumber at a time.}

\begin{example}[Rich and poor solutions of the heat equation]\label{ex:heat-rich-poor}
A single heat mode $e^{-\nu\kappa^2t+i\kappa x}$ and any finite sum of such modes belong to $\EP$.
The heat kernel $t^{-1/2}e^{-x^2/(4\nu t)}$ does not: its $x$-derivatives are polynomials of
increasing degree times the same Gaussian, so they span an infinite-dimensional space.
\end{example}
